\section*{Abstract}
\addcontentsline{toc}{section}{Abstract}

VIGÍA integrates publicly available hazard detectors under a single gated
pipeline and adds a shared temporal validation layer that makes them usable on
ordinary municipal cameras. The system is explicitly an integration project:
every detector it runs was trained by someone else and is credited by name,
licence and source. The contribution is the validation layer and the
evaluation methodology that measures what that layer buys and what it costs.

The problem is not detection accuracy. Published hazard detectors are good and
freely available; what prevents their deployment is that they alarm far too
often. Measured on its authors' own validation split, the PyroNear wildfire
detector raises an alarm on \FireImageFAR{} of frames containing no smoke
(95\,\% interval \FireImageFARCI, \FireNegatives{} negatives). That is not a
defect but the operating point its authors chose, on the assumption that
something downstream filters the output. VIGÍA is that something.

Across \AblationNegativeSequences{} negative sequences the validator reduces
the false-alarm rate from \RawFAR{} to \ValidatedFAR{}, a \FARReduction\,\%
reduction, while still detecting every one of the \FiresTotal{} ignition
sequences, at a cost of \AddedMinutesToDetect{} minutes of additional latency
which this document derives in closed form instead of leaving as an observation. We then
run the control most system papers omit — matching that false-alarm rate by
raising the detector's confidence floor — and report that on this negative set
an oracle-tuned floor achieves the same rate at higher recall. A
cross-camera transfer experiment resolves the objection this raises: a floor
tuned on one set of cameras and applied to unseen cameras either lets the
false-alarm rate drift to \TransferTestFAR{} or loses a fire (recall
\TransferBARecall), while the cascade's untuned persistence holds both
false-alarm rate and recall across the same held-out cameras. The validator's
defensible value is thus that its stability does not depend on access to the
deployment's negative distribution, and that it deduplicates alerts in a way
no threshold can (\cref{sec:threshold-matched}). The same validator, changed
only by configuration, transfers across hazards whose detectors were built by
different people, on different data, for phenomena with different physics; a
build-failing test enforces that no hazard-specific branch exists in it.

Two commitments are enforced in code, not stated in prose. The runtime
imports no AGPL-licensed package, and no biometric identification code path
exists; both are checked by guards that parse the source tree and fail the
build. A third guard checks every published figure against the results file it
cites.

Where the evidence is thin, this document says so and quantifies it. Building
access rests on \BuildingTestImages{} held-out images: its image-level recall
is \BuildingImageRecall, but the interval that sample supports is
\BuildingImageRecallCI, and the paper reports the second.

\medskip
\noindent\textbf{Keywords:} wildfire smoke detection; flood monitoring;
multi-hazard early warning; temporal validation; false-alarm reduction;
camera networks; disaster risk reduction; ONNX Runtime.
